\section{Further properties of projections}

1 Also, what would the analogies with characteristic functions look like
2 Is this also true for characteristic functions? 
4 Is this also true for characteristic functions?
5
6
7
8 Again, make the link between characteristic functions. What should perpendicular characteristic functions translate to?
9
10 So reduction can be stated in terms of projection operators, making projection operators particularly powerful. Does this make sense for characteristic functions too?
11
10
14
15 This is a crucial step in proving the finite-dimensional version of the spectral theorem. The idea is to first find an eigenvector, and pick Y = the span of this eigenvector. Then T splits up into two smaller-dimensional operators, on which we can induct!

\begin{exercise}{1}
fill
\end{exercise}
\begin{proof}
fill
\end{proof} 

\begin{exercise}{2}
fill
\end{exercise}
\begin{proof}
fill
\end{proof} 

\begin{exercise}{4}
fill
\end{exercise}
\begin{proof}
fill
\end{proof} 

\begin{exercise}{5}
fill
\end{exercise}
\begin{proof}
fill
\end{proof} 

\begin{exercise}{6}
fill
\end{exercise}
\begin{proof}
fill
\end{proof} 

\begin{exercise}{7}
fill
\end{exercise}
\begin{proof}
fill
\end{proof} 

\begin{exercise}{8}
fill
\end{exercise}
\begin{proof}
fill
\end{proof} 

\begin{exercise}{9}
fill
\end{exercise}
\begin{proof}
fill
\end{proof} 

\begin{exercise}{10}
fill
\end{exercise}
\begin{proof}
fill
\end{proof} 

\begin{exercise}{11}
fill
\end{exercise}
\begin{proof}
fill
\end{proof} 

\begin{exercise}{12}
fill
\end{exercise}
\begin{proof}
fill
\end{proof} 

\begin{exercise}{13}
fill
\end{exercise}
\begin{proof}
fill
\end{proof} 

\begin{exercise}{14}
fill
\end{exercise}
\begin{proof}
fill
\end{proof} 

\begin{exercise}{15}
fill
\end{exercise}
\begin{proof}
fill
\end{proof} 
